\ifdefined\standalone 
\documentclass[12pt]{article}
\usepackage{graphicx}
\usepackage{float}
\usepackage[a4paper, margin=1in]{geometry}
\usepackage[utf8]{inputenc}
\usepackage{textcomp}
\usepackage{amsmath}
\usepackage{booktabs}
\usepackage{pgfplotstable}
\usepackage{booktabs}
\usepackage[hidelinks]{hyperref} % <- hypertext
\providecommand{\SPATIALDIR}{.}

\begin{document}
\fi

\section{Spatial Convergence Verification}

This verification case assesses the spatial convergence behavior of the \texttt{Gpyro} solver in a cone calorimeter configuration. All simulations are run using a fixed time step \(\Delta t = 0.05\,\mathrm{s}\), while the spatial grid size \(\Delta z\) is varied from \(1\,\mathrm{mm}\) to \(0.0001\,\mathrm{mm}\).

Figure~\ref{fig:mlr_curves_dz} presents the mass loss rate obtained for different spatial discretizations, and Table~\ref{tab:cpu_vs_error_space} summarizes the corresponding mean errors and CPU times.

A log-log plot of the error versus \(\Delta z\) is shown in Figure~\ref{fig:error_convergence_space}. The linear regression indicates that the numerical scheme exhibits a spatial convergence rate close to \(\mathcal{O}(\Delta z^1)\).

\begin{figure}[H]
\centering
\begin{minipage}[t]{0.48\textwidth}
\centering
\includegraphics[width=\linewidth]{grid_convergence.png}
\caption{Mass loss rate profiles for different grid sizes.}
\label{fig:mlr_curves_dz}
\end{minipage}%
\begin{minipage}[t]{0.48\textwidth}
\centering
\includegraphics[width=\linewidth]{grid_error.png}
\caption{Mean error vs. grid size \(\Delta z\).}
\label{fig:error_convergence_space}
\end{minipage}
\hfill
\end{figure}

\begin{table}[H]
\centering
\caption{Mean error and CPU time for different spatial grid sizes.}
\label{tab:cpu_vs_error_space}
\pgfplotstableread[col sep=comma]{\SPATIALDIR/convergence_summary.csv}\datatable
\pgfplotstabletypeset[
    col sep=comma,
    string type,
    columns={dz (mm),Mean Error,CPU Time (s)},
    columns/dz (mm)/.style={column name={\(\Delta z\) (mm)}},
    columns/Mean Error/.style={column name={Mean Error}},
    columns/CPU Time (s)/.style={column name={CPU Time (s)}},
    every head row/.style={before row=\toprule, after row=\midrule},
    every last row/.style={after row=\bottomrule}
]{\datatable}
\end{table}

\ifdefined\standalone
\end{document}
\fi



